\section{相同执行状态下的恢复评价设计}
\label{sec:same-state-design}

自由运行准确率刻画模型从问题开始的完整求解表现。同状态评价进一步控制模型接手的历史，以研究普通续训和恢复续训在相同正常及偏离状态下的继续求解能力。该实验已固定2400题候选池、状态构造、检查点和统计规则，实际模型续行将在计算环境与原检查点核验后实施。

由固定初始模型I采集候选问题的执行轨迹，按失败结果、共同成功前缀、参考后缀可执行性及预算条件确定合格样本。设第$q$题共同前缀为$H_q$，实际分歧动作及观察为$d_q$。正常条件$N$使用$H_q$，偏离条件$D$使用$H_q\Vert d_q$；每个条件内，A与C读取相同历史，并重放恢复相同内部句柄状态。

\begin{table}[htbp]
\centering\small
\caption{同状态评价的模型与历史条件}
\label{tab:same-state-factorial}
\begin{tabular}{lll}
\toprule
模型 & 正常历史$N$ & 偏离历史$D$\\
\midrule
普通续训A & A-N：共同成功前缀 & A-D：前缀及实际分歧\\
恢复续训C & C-N：同一成功前缀 & C-D：同一前缀及分歧\\
\bottomrule
\end{tabular}
\end{table}

两条件均按$24-k_q-1$次剩余调用和2048个新增生成token设置续行预算，其中$k_q$为共同前缀的调用数。每步生成上限192 token，上下文上限8192 token；偏离历史增加的上下文长度与累计输入开销分别记录。候选筛选基于I的轨迹和参考重放，在读取A/C结果前完成，全部排除原因保留。

设$Y_{q,s}^{M,h}\in\{0,1\}$表示种子$s$、模型$M$在历史条件$h$下的最终正确性，以两个既定续训种子的题级平均差构造
\begin{equation}
 \Delta_h=\frac{1}{n}\sum_{q=1}^{n}\frac{1}{2}\sum_{s=1}^{2}
 \left(Y_{q,s}^{C,h}-Y_{q,s}^{A,h}\right),\qquad h\in\{N,D\}.
\end{equation}
预登记主指标为$\Delta_D$，同时报告$\Delta_N$及交互差异$\Delta_D-\Delta_N$。配对bootstrap以问题为抽样单位，同时保留其两个种子与两种历史条件。偏离条件的正向增量刻画该类状态中的求解收益，交互差异用于分析收益在两类历史中的分布；结果结合完成率、无效调用、来源状态分层和成本共同解释。
